\scanpage{199}
\mspage{73}
\uline{Corollaire} \uline{5.16} (de prop 5.13.) Toute réunion finie de parties L-accessibles de $M$ est L-accessible. \unsure{Toute} intersection $\bigcap_{\alpha\in A}M_\alpha$ de \inserted{une famille \unsure{petite} de} parties L-accessibles de $M$, est L-accessible.

6. \unsure{Essay|Essai} avorté pour un critère d'engendrement des parties L-accessibles de $\unsure{\mathrm{Fl}|F\ell}\,M$.

Soit $M$ une catégorie munie d'une filtration cardinale \inserted{$(\mathrm{Filt}^{\unsure{\pi|R}})_{\pi\geq\pi_0}$}, et \deleted{\unsure{considérons}} soit $S$ un objet de $\mathrm{Cat}$, considérons
\[
 M(S)=\underline{\mathrm{Hom}}\,(S,M),
\]
et munissons $M(S)$ de la filtration par les
\[
 M^\pi(S)=\underline{\mathrm{Hom}}\,\bigl(S,M^\pi=\underline{\mathrm{Filt}}^\pi(M)\bigr).
\]
Je dis que c'est là une filtration cardinale de $M(S)$. \deleted{\unsure{Il}} \unsure{faut} vérifier les conditions a) b) c) de SGA 4 I 9.12. Les conditions a) ($\mathrm{Filt}^\pi$ \unsure{ess. petite|ens. petits.}) et b (condition $\mathrm{L}_{\pi_0}$, et stabilité \unsure{de} $\mathrm{Filt}^\pi$ par $\underset I{\underline{\lim}}$, \deleted{\unsure{filtran}} $I$ filtrant grand devant $\pi_0$ et de cardinal $\leq\pi$) passent aussitôt de $M$ à $M(S)$. La condition c) est plus délicate. (\unsure{représentation}
% archive sheet number: 198
